%====================================================================
%   SECTION 5 — BACK-TEST PROTOCOL AND METRICS
%====================================================================

\section{Back-test protocol and metrics}\label{sec:backtest}

\paragraph{Costs and walk-forward.}\label{sec:walkforward}
Every position change in leg $i$ pays the half-spread $\mathrm{HS}_t^{\,i}$ of \eqref{eq:midprice} plus a fixed commission \parencite{aldridge2013hft}, so the round-trip cost of one pair-position is approximately $\mathrm{HS}^{A} + \hat\beta\,\mathrm{HS}^{B}$ in basis points. Strategies are evaluated walk-forward: parameters $\Theta$ and hyperparameters $(W_\beta, W_z, z_q, z_x, z_v, \delta)$ are fitted on a training segment and applied to a disjoint out-of-sample test segment, with windows slid forward \parencite{lopezdeprado2018advances}.

\paragraph{Core performance metrics.}\label{sec:metrics}
We report three core metrics in the body of the thesis: annualised volatility, the Sharpe ratio, and the maximum drawdown. A fuller set of distributional, tail-risk, and risk-adjusted diagnostics (skewness, excess kurtosis, VaR, CVaR, Sortino, Calmar, Ulcer index) is collected in Appendix~\ref{sec:apx-extra-metrics} and reported in the per-pair tearsheets. Let $R_t$ denote the strategy return at bar $t$, $\mu_R = \mathbb{E}[R]$, $\sigma_R = \sqrt{\mathrm{Var}(R)}$, $V_t$ the portfolio value with running peak $V_t^{\star} = \max_{s \leq t} V_s$, and $R_f$ a benchmark risk-free rate (set to zero). The three core metrics are
\begin{equation}\label{eq:coremetrics}
    \sigma_{\mathrm{ann}} \;=\; \sigma_R\sqrt{T},
    \qquad
    \mathrm{SR} \;=\; \frac{\mu_R - R_f}{\sigma_R},
    \qquad
    \mathrm{MDD} \;=\; \max_{t}\,\frac{V_t^{\star} - V_t}{V_t^{\star}},
\end{equation}
covering scale (volatility), risk-adjusted return (Sharpe \parencite{sharpe1994ratio}), and worst-case capital loss (maximum drawdown). Trade-level diagnostics (win rate, market exposure, profit factor, payoff ratio) summarise frequency and size; see Section~\ref{sec:result} for their interpretation.
